%% Définition géométrique du produit scalaire : deux vecteurs de même origine et l'angle
%% qu'ils forment. Le signe du produit est celui de cos(theta) (ici angle aigu : produit positif).
\documentclass[tikz,border=4pt]{standalone}
\usepackage{liaisons}
\begin{document}
\begin{tikzpicture}[x=1.6cm,y=1.6cm,
  vect/.style={-{Stealth[length=7pt]}, line width=1.4pt},
  arc/.style={line width=0.8pt, black!70},
  long/.style={font=\small, text=black!55, inner sep=3pt}]

\coordinate (O) at (0,0);
\coordinate (U) at (12:3.2);    % u : long, presque horizontal
\coordinate (V) at (62:2.1);    % v : plus court, incliné d'environ 50° au-dessus de u

%% ---------------- les deux vecteurs -------------------------------------
\draw[vect, solideF] (O) -- (U)
      node[midway, sloped, below, long] {$\left\|\vec{u}\right\|$}
      node[anchor=west, inner sep=3pt, font=\small, text=solideF] {$\vec{u}$};
\draw[vect, solideB] (O) -- (V)
      node[midway, sloped, above, long] {$\left\|\vec{v}\right\|$}
      node[anchor=south, inner sep=3pt, font=\small, text=solideB] {$\vec{v}$};
\fill (O) circle (1.7pt);
\node[font=\small, anchor=north east, inner sep=2pt] at (O) {$O$};

%% ---------------- angle theta entre les deux vecteurs --------------------
\draw[arc] (12:0.85) arc (12:62:0.85);
\node[font=\small, anchor=west, inner sep=2pt] at (37:0.95) {$\theta$};

%% ---------------- l'égalité ----------------------------------------------
\node[font=\small, anchor=north, inner sep=2pt] at (1.7,-0.6)
     {$\vec{u}\cdot\vec{v}=\left\|\vec{u}\right\|\times\left\|\vec{v}\right\|\times\cos(\theta)$};
\end{tikzpicture}
\end{document}
